\section{EA Test 3: momentum conservation over a small obstruction}

Test~3 of the Environment Agency benchmark \cite{NeelzPender2013} asks
whether a model carries momentum over a low obstruction. A
300~m~$\times$~100~m channel falls at 1:200 into two depressions, centred
on points~1 ($x = 150$~m) and~2 ($x = 250$~m), separated by a crest at
$x = 200$~m (Figure~\ref{fig:ea3_setup}). A short inflow of 1,310~m$^3$
across the upstream end just fills the first depression. A model that
conserves momentum carries some of the water over the crest into the
second; one that does not leaves it all in the first. Manning's
$n = 0.01$, the specified resolution is 5~m (here triangles of at most
25~m$^2$, about 1,200 in all), the walls are closed and the run
lasts 900~s.

\begin{figure}[h]
\begin{center}
\includegraphics[width=0.95\textwidth]{setup.png}
\end{center}
\caption{EA Test 3: the bed along the channel centreline and the inflow.}
\label{fig:ea3_setup}
\end{figure}

The report gives no tabulated results. In its summary of the full
shallow-water packages (Section~4.4.3), the level at point~1 rises from
the bed at 9.75~m to about 9.98--9.99~m by $t \approx 120$~s, and from
$t \approx 150$~s water crosses the crest and the level at point~2 rises by
about 5--6~cm. These are the grey bands of Figure~\ref{fig:ea3_stage}, and
with the volume balance and a stored baseline per flow algorithm they are
the checks of Figure~\ref{fig:ea3_checks}.

\begin{figure}
\begin{center}
\includegraphics[width=0.95\textwidth]{stage_plot.png}
\end{center}
\caption{EA Test 3: water level at points 1 and 2, with the report's
summary of the full shallow-water packages.}
\label{fig:ea3_stage}
\end{figure}

\begin{figure}
\begin{center}
\includegraphics[width=0.95\textwidth]{speed_plot.png}
\end{center}
\caption{EA Test 3: speed at points 1 and 2.}
\end{figure}

\begin{figure}
\begin{center}
\includegraphics[width=0.9\textwidth]{checks.png}
\end{center}
\caption{EA Test 3: the checks and their outcome.}
\label{fig:ea3_checks}
\end{figure}
